\section[nosectionpage]{Ordinary \texttt{SocketServer} objects and the MQTTbroker application}

\begin{frame}[fragile]{Ordinary \texttt{SocketServer} objects and the MQTTbroker application}
  \begin{itemize}
    \item The MQTTbroker \texttt{\alert{mqttbroker}} is an ordinary \alert{TCP server} application communicating \alert{unencrypted} via \alert{IPv4} with the MQTT broker protocol attached.
    Thus you have to use one of the \alert{template classes \texttt{SocketServer}} provided by SNode.C and provide a \texttt{\alert{SocketContextFactory}} class as \alert{template argument}.

    \begin{cppcode}[gobble=4]
      #include <net/in/stream/legacy/SocketServer.h> // stream => TCP (dgram would be UDP)
      #include "SharedSocketContextFactory.h"

      int main(int argc, char* argv[]) {
        ...
        net::in::stream::legacy::SocketServer<mqtt::mqttbroker::SharedSocketContextFactory> server(...);
        ...
      }
    \end{cppcode}
    \item The \texttt{\alert{SocketContextFactory}} is responsible for \alert{creating} a concrete \texttt{\alert{SocketContext}} object for \alert{every connected client}, \alert{representing} the \alert{application protocol} (e.g. \texttt{\alert{HTTP}}, or \texttt{\alert{Websocket}}, or \texttt{\alert{MQTT}}, or \ldots) used by the application.
  \end{itemize}
\end{frame}
